\documentclass[12pt,letterpaper]{article}
\usepackage[utf8]{inputenc}
\usepackage[letterpaper, margin=1in, bottom=2.5in, top=2.2in]{geometry}
\usepackage{fancyhdr}
\setlength{\headheight}{15pt}
\setlength{\footskip}{72pt}

\usepackage{xcolor}
\usepackage{amsmath}
\usepackage{anyfontsize}
\usepackage{graphicx}
\usepackage{tikz}
\usepackage{listings}
\usepackage{sourcesanspro}
\renewcommand{\familydefault}{\sfdefault}
\usepackage{fontawesome5}
\usepackage{titlesec}
\usepackage{setspace}
\usepackage{hyperref}
\usepackage{mdframed}
\usepackage{qrcode}
\usetikzlibrary{calc,shapes,positioning}
\usepackage{eso-pic}

% Command for inline code highlighting
\newcommand{\code}[1]{\texttt{\textcolor{accentColor}{#1}}}

\AtBeginDocument{\color{primaryColor}}

% Colors for Python theme
\definecolor{bgColor}{RGB}{13, 17, 23}  % Dark background
\definecolor{primaryColor}{RGB}{255, 255, 255}  % White text
\definecolor{accentColor}{RGB}{68, 171, 255}  % Python blue for main accent
\definecolor{accentColor2}{RGB}{200, 160, 35}  % Mustard yellow for highlighting concepts
\definecolor{pythonSecondary}{RGB}{180, 180, 180}  % Light gray
\definecolor{secondaryColor}{RGB}{230, 235, 240}  % Almost white gray
\definecolor{terminalBg}{RGB}{22, 22, 22}  % Terminal background
\definecolor{terminalFrame}{RGB}{40, 40, 40}  % Terminal frame
\definecolor{lineNumberColor}{RGB}{100, 100, 100}  % Line number color
\definecolor{dividerColor}{RGB}{120, 130, 150}  % Divider color

% Code colors
\definecolor{codeTextColor}{RGB}{255, 255, 255}  % Basic text
\definecolor{codeKeywordColor}{RGB}{68, 171, 255}  % Keywords in Python blue
\definecolor{codeCommentColor}{RGB}{98, 150, 85}  % Comments 
\definecolor{codeStringColor}{RGB}{200, 160, 35}  % Strings in mustard yellow
\definecolor{codeClassColor}{RGB}{78, 201, 176}  % Classes
\definecolor{codeFunctionColor}{RGB}{200, 160, 35}  % Functions in mustard yellow
\definecolor{codeNumberColor}{RGB}{181, 206, 168}  % Numbers

\pagecolor{bgColor}

\hypersetup{
    colorlinks=true,
    linkcolor=accentColor,
    filecolor=accentColor,
    urlcolor=accentColor,
}

% Define Python language for listings
\lstdefinelanguage{Python}{
  keywords={and, as, assert, break, class, continue, def, del, elif, else, except, finally, for, from, global, if, import, in, is, lambda, nonlocal, not, or, pass, raise, return, try, while, with, yield, async, await},
  sensitive=true,
  comment=[l]{\#},
  morestring=[b]',
  morestring=[b]"
}

% Listings configuration
\lstset{
  language=Python,
  basicstyle=\ttfamily\bfseries\color{codeTextColor},
  backgroundcolor=\color{terminalBg},
  commentstyle=\color{codeCommentColor},
  keywordstyle=\color{codeKeywordColor},
  stringstyle=\color{codeStringColor},
  numberstyle=\color{lineNumberColor},
  breaklines=true,
  breakatwhitespace=true,
  tabsize=4,
  showstringspaces=false,
  frame=none,
  xleftmargin=15pt,
  xrightmargin=0pt,
  aboveskip=5pt,
  belowskip=5pt,
  numbers=left,
  numbersep=8pt,
  extendedchars=true,
  keepspaces=true,
  columns=flexible,
  lineskip=2pt,
  emph={[2]print,len,type,id,dict,list,tuple,set,str,int,float,bool,None,True,False},
  emphstyle={[2]\color{codeFunctionColor}}
}

\newenvironment{macterminal}{%
    \begin{mdframed}[
        linecolor=terminalFrame,
        backgroundcolor=terminalBg,
        roundcorner=5pt,
        skipabove=5pt,
        skipbelow=5pt,
        linewidth=1pt,
        innertopmargin=5pt,
        frametitle={%
            \tikz[baseline=(current bounding box.east), outer sep=0pt]{
                \fill[red!80!black] (0,0) circle (5pt);
                \fill[yellow!80!black] (0.7,0) circle (5pt);
                \fill[green!70!black] (1.4,0) circle (5pt);
            }
        },
        frametitlealignment=\raggedright,
        frametitleaboveskip=8pt,
        frametitlebelowskip=0pt,
    ]
}{%
    \end{mdframed}%
}

\newcommand{\verspace}{\vspace{5pt}}
\newcommand{\languagetag}[1]{
    \begin{tikzpicture}[baseline]
        \node[fill=accentColor, text=primaryColor, rounded corners=5pt, inner sep=7pt] {
            {\normalsize\textbf{#1}}
        };
    \end{tikzpicture}
}

% Title formatting
\titleformat{\section}
  {\LARGE\bfseries\color{primaryColor}}
  {\thesection. }
  {0pt}
  {}
  []

\titleformat{\subsection}
  {\Large\bfseries\color{accentColor}}
  {\thesubsection. }
  {0pt}
  {}
  []

\titleformat{\subsubsection}
  {\large\bfseries\color{accentColor2}}
  {\thesubsubsection. }
  {0pt}
  {}
  []

\titlespacing*{\section}{0pt}{18pt}{8pt}
\titlespacing*{\subsection}{0pt}{12pt}{5pt}
\titlespacing*{\subsubsection}{0pt}{8pt}{3pt}

% Header and footer setup
\pagestyle{fancy}
\fancyhf{}

\fancyhead[L]{
    \begin{tikzpicture}[remember picture, overlay]
        \fill[accentColor, rounded corners=3pt] (0,0) rectangle (2.2cm,0.7cm);
        \node[text=primaryColor, font=\bfseries] at (1.1cm,0.35cm) {PYTHON};
        \fill[primaryColor] (2.35cm,0.1cm) rectangle (2.40cm,0.6cm);
        \node[text=primaryColor, font=\bfseries, anchor=west] at (2.35cm,0.35cm) {NAMESPACES};
    \end{tikzpicture}
}

\renewcommand{\headrulewidth}{0pt}
\renewcommand{\footrulewidth}{0.2pt}
\renewcommand{\footruleskip}{1cm}

\fancyfoot[C]{
    \vspace*{0.1cm}
    \noindent
    \begin{minipage}{\textwidth}
        \begin{flushleft}
            % Profile photo
            \raisebox{0.7cm}{
            \begin{tikzpicture}[baseline]
                \path[fill=bgColor] (0,0) circle (0.6cm);
                \clip (0,0) circle (0.6cm);
                \node at (0,0) {
                    \includegraphics[width=1.2cm,height=1.2cm]{../../Images/profile-image.jpeg}
                };
            \end{tikzpicture}
            }
            % Profile information
            \begin{minipage}[b]{0.5\textwidth}
                {\large\bfseries\color{primaryColor}Alejandro Sánchez Yalí}
                \par\vspace{1pt}
                {\small\color{secondaryColor}Software Developer | AI \& Blockchain Enthusiast}
                \par\vspace{1pt}
                {\small\color{accentColor2}\faGlobe\hspace{5pt}\color{secondaryColor}www.asanchezyali.com}
            \end{minipage}
        \end{flushleft}
        \vspace{6pt}
    \end{minipage}
}

\fancypagestyle{plain}{
    \fancyhf{}
    \renewcommand{\headrulewidth}{0pt}
    \renewcommand{\footrulewidth}{0pt}
}

\renewcommand{\labelitemi}{\textcolor{accentColor}{$\bullet$}}
\renewcommand{\labelitemii}{\textcolor{accentColor2}{$\circ$}}

\usepackage{relsize}
\AtBeginDocument{\relsize{1}}

% Command for elegant QR
\newcommand{\elegantqr}[2]{
    \qrcode[height=2.5cm]{#1}
    \\[0.1cm]
    {\hspace{0.2cm}\color{primaryColor}\small #2\par}
}


\usepackage[utf8]{inputenc}
\lstset{inputencoding=utf8}

\newcommand{\BackgroundPic}{%
    \put(0,0){%
        \begin{tikzpicture}[remember picture, overlay]
            \node[inner sep=0pt] at (current page.center) {
                \includegraphics[width=\paperwidth,height=\paperheight,keepaspectratio=false]{../../Images/Futuristic Cybersecurity Analyst}
            };
            \fill[black, opacity=0.7] (current page.south west) rectangle (current page.north east);
        \end{tikzpicture}%
    }%
}

\newcommand{\titlepagecontents}{%
    \AddToShipoutPicture*{\BackgroundPic}
    \vspace*{2cm}
    \begin{flushleft}
    \languagetag{Python}\\[0.4cm]
    {\fontsize{38}{52}\bfseries\color{primaryColor}Mastering \color{accentColor}Python\\\color{accentColor}Regular Expressions\par}
    \vspace{0.3cm}
    {\fontsize{18}{52}\color{secondaryColor}Pattern Matching Made Simple\par}
    \vspace{0.3cm}
    {\color{secondaryColor}\today\par}
    \vspace{1.8cm}
    \elegantqr{https://github.com/asanchezyali/social-media-posts/tree/main/Python/Regex}{Source Code}
    \end{flushleft}
}

\newcommand{\finalpagecontents}{%
    \AddToShipoutPicture*{\BackgroundPic}
    \vspace*{3cm}
    \begin{flushleft}
    \languagetag{Feedback}\\[0.4cm]
    {\fontsize{46}{52}\bfseries\color{primaryColor}Found this \color{accentColor}helpful?\par}
    \vspace{0.3cm}
    {\fontsize{18}{52}\color{secondaryColor}Save, comment and share\par}  
    \vspace{0.3cm}
    {\color{secondaryColor}\today\par}
    \end{flushleft}
}

% Update header for Regex
\fancyhead[L]{
    \begin{tikzpicture}[remember picture, overlay]
        \fill[accentColor, rounded corners=3pt] (0,0) rectangle (2.2cm,0.7cm);
        \node[text=primaryColor, font=\bfseries] at (1.1cm,0.35cm) {PYTHON};
        \fill[primaryColor] (2.35cm,0.1cm) rectangle (2.40cm,0.6cm);
        \node[text=primaryColor, font=\bfseries, anchor=west] at (2.35cm,0.35cm) {REGEX};
    \end{tikzpicture}
}

\begin{document}
\setlength{\parindent}{0pt}
\begin{titlepage}
    \titlepagecontents
\end{titlepage}

\section{Introduction to Regular Expressions}

Regular expressions (regex) are powerful tools for pattern matching and text manipulation in Python. They provide a concise and flexible way to:

\begin{itemize}
    \item Search and match specific patterns in text
    \item Validate string formats (emails, phone numbers, etc.)
    \item Extract information from structured text
    \item Perform complex string replacements
\end{itemize}

Python's \code{re} module provides comprehensive support for regular expressions, making it an essential tool for text processing tasks.

% Quick Reference Table for Regex Syntax

\section*{Regex Quick Reference}

\begin{tabular}{ll}
\textbf{Pattern} & \textbf{Meaning} \\
\hline
\texttt{.} & Any character except newline \\
\texttt{*} & Zero or more repetitions \\
\texttt{+} & One or more repetitions \\
\texttt{?} & Zero or one repetition \\
\texttt{[a-z]} & Any lowercase letter \\
\texttt{[A-Z]} & Any uppercase letter \\
\texttt{[0-9]} & Any digit \\
\texttt{[abc]} & a, b, or c \\
\texttt{\textbackslash d} & Digit (0-9) \\
\texttt{\textbackslash w} & Word character (a-z, A-Z, 0-9, \_) \\
\texttt{\textbackslash s} & Whitespace \\
\texttt{\^} & Start of string \\
\texttt{\$} & End of string \\
\texttt{|} & Alternation (or) \\
\texttt{()} & Grouping \\
\texttt{\textbackslash} & Escape special character \\
\end{tabular}

\section{Basic Pattern Matching}

Let's start with fundamental regex concepts and pattern matching:

\begin{macterminal}
\lstinputlisting{codes/01_basic_patterns.py}
\end{macterminal}

Key concepts demonstrated:
\begin{itemize}
    \item The \code{r} prefix creates a raw string
    \item \code{re.search()} finds the first match in text
    \item \code{re.IGNORECASE} flag makes the search case-insensitive
    \item \code{re.findall()} returns all matches
    \item Word boundaries with \code{\textbackslash b}
\end{itemize}

\section{Visual Examples}

\textbf{Example 1: Simple Match}

\begin{verbatim}
Pattern: fox
String:  The quick brown fox jumps over the lazy dog
Match:   fox
\end{verbatim}

\textbf{Example 2: Digit Extraction}

\begin{verbatim}
Pattern: \d+
String:  Order numbers: 123, 456, 789
Match:   123, 456, 789
\end{verbatim}

\section{Common Pattern Validation}

Regular expressions are frequently used to validate common text formats. Here are practical examples:

\begin{macterminal}
\lstinputlisting{codes/02_pattern_validation.py}
\end{macterminal}

Important pattern components:
\begin{itemize}
    \item \code{\string^} and \code{\string$} anchor the pattern to start/end
    \item \code{[a-zA-Z0-9]} matches any alphanumeric character
    \item \code{+} matches one or more occurrences
    \item \code{*} matches zero or more occurrences
    \item \code{?} makes a character optional
    \item \code{\textbackslash d} matches any digit
    \item \code{\textbackslash w} matches word characters
\end{itemize}

\section{Advanced Features}

Regular expressions support advanced features such as:

\begin{itemize}
    \item \textbf{Named capture groups:} \code{(?P<name>pattern)}
    \item \textbf{Non-capturing groups:} \code{(?:pattern)}
    \item \textbf{Positive lookahead:} \code{(?=pattern)}
    \item \textbf{Negative lookahead:} \code{(?!pattern)}
    \item \textbf{Positive lookbehind:} \code{(?<=pattern)}
    \item \textbf{Negative lookbehind:} \code{(?<!pattern)}
    \item \textbf{Greedy vs. non-greedy quantifiers:} \code{*} vs. \code{*?}
    \item \textbf{Backreferences:} \code{\textbackslash 1}, \code{\textbackslash 2}, etc.
    \item \textbf{Pattern substitution:} \code{re.sub()}
\end{itemize}

\textbf{Visual Example: Non-greedy Matching}

\begin{verbatim}
Pattern: <.*?>
String:  <tag>content</tag>
Match:   <tag>, </tag>
\end{verbatim}

\section{Testing Your Regex}

You can experiment with regular expressions interactively using online tools such as:
\begin{itemize}
    \item \href{https://regex101.com/}{regex101.com}
    \item \href{https://pythex.org/}{pythex.org}
    \item \href{https://regexr.com/}{regexr.com}
\end{itemize}

\section{Common Regex Recipes}

\begin{itemize}
    \item \textbf{Extract emails:} \verb|[a-zA-Z0-9._%+-]+@[a-zA-Z0-9.-]+\.[a-zA-Z]{2,}|
    \item \textbf{Extract URLs:} \verb|https?://[\w.-]+(?:/[\w./?%&=-]*)?|
    \item \textbf{Extract dates (YYYY-MM-DD):} \verb|\d{4}-\d{2}-\d{2}|
    \item \textbf{Extract phone numbers:} \verb|\+?1?\d{9,15}|
\end{itemize}

\section{Error Handling and Useful Flags}

\begin{itemize}
    \item Use \code{try/except} to catch \code{re.error} exceptions in Python.
    \item Useful flags:
    \begin{itemize}
        \item \code{re.IGNORECASE} (case-insensitive)
        \item \code{re.MULTILINE} (\verb|^| and \verb|$| match at line breaks)
        \item \code{re.DOTALL} (dot matches newline)
        \item \code{re.UNICODE} (Unicode matching)
        \item \code{re.VERBOSE} (allow comments and whitespace in pattern)
    \end{itemize}
\end{itemize}

\section{Best Practices}

When working with regular expressions in Python:

\begin{itemize}
    \item Use raw strings (\code{r"pattern"}) to avoid escaping backslashes
    \item Start with simple patterns and gradually add complexity
    \item Test patterns with a variety of inputs
    \item Use \code{re.VERBOSE} for complex patterns with comments
    \item Consider performance with large texts
    \item Document complex patterns
    \item Use existing patterns for common formats
\end{itemize}

\section{Common Pitfalls}

Be aware of these common regex pitfalls:

\begin{itemize}
    \item Catastrophic backtracking with nested quantifiers
    \item Greedy vs. non-greedy matching
    \item Unicode handling in patterns
    \item Over-complicated patterns
    \item Not escaping special characters
    \item Incorrect character class usage
\end{itemize}

\section{Performance Tips}

Optimize your regex usage:

\begin{itemize}
    \item Compile patterns you use repeatedly with \code{re.compile()}
    \item Use non-capturing groups \code{(?:\string.\string.\string.)} when possible
    \item Avoid unnecessary backtracking
    \item Be specific with character classes
    \item Consider alternative solutions for complex patterns
    \item Profile pattern matching on large datasets
\end{itemize}

\section{References}

\begin{itemize}
    \item Python Documentation. (2025). \textit{re - Regular expression operations}.
    \href{https://docs.python.org/3/library/re.html}{Link}
    
    \item Goyvaerts, J., \& Levithan, S. (2023). \textit{Regular Expressions Cookbook}. O'Reilly Media.
    
    \item This article was edited and written in collaboration with AI. If you find any inconsistencies or have suggestions for improvement, please don't hesitate to open an issue in our \href{https://github.com/asanchezyali/social-media-posts}{GitHub} repository or reach out directly.
\end{itemize}

\section{Explore More Python Concepts}
\vspace{10pt}
\begin{tikzpicture}
    \fill[color=terminalBg, rounded corners=5pt] (0,0) rectangle (\textwidth,-5);
    
    \pgfmathsetmacro{\availableWidth}{\textwidth-4cm}
    
    \begin{scope}
        \node[text=accentColor, font=\Large\bfseries, anchor=north west] at (0.5cm,-0.5cm) {Enjoyed This Content?};
        
        \node[text=primaryColor, font=\normalsize, anchor=north west] at (0.5cm,-1.2cm) {Don't miss my other Python articles:};
        
        \node[text=secondaryColor, font=\large\bfseries, anchor=north west, text width=\availableWidth] at (0.5cm,-1.9cm)
            {\color{accentColor}Python Namespaces: \color{primaryColor}A Deep Dive into Name Resolution};
        
        \node[text=pythonSecondary, font=\normalsize, text width=\availableWidth, anchor=north west] at (0.5cm,-3.1cm) {
            Learn how Python organizes and manages variable names, scopes, and namespaces.
        };
    \end{scope}
    
    \node[anchor=center] at ({\textwidth-1.5cm}, {-2.5cm}) {
        \qrcode[height=2.2cm]{https://www.linkedin.com/feed/update/urn:li:activity:7324383306865664000?utm_source=share&utm_medium=member_desktop&rcm=ACoAACfC_XIB4R901-Rs4Cep7e9KjyGl5ZhS47o}
    };
    
\end{tikzpicture}
\vspace{5pt}

\clearpage
\thispagestyle{empty}
\finalpagecontents

\end{document}